\section{Experiments}
\label{sec:experiments}

We evaluate 14 algorithms across 12 synthetic configurations and 2 real datasets, examining performance across regimes defined by topic count~$K$ and decay rate~$\lambda$.

\subsection{Experimental Setup}
\label{sec:setup}

\paragraph{Synthetic environment.}
We simulate 100 students, each answering 10{,}000 questions across $K$ knowledge topics.
We vary $K \in \{6, 20, 50, 100\}$ and $\lambda \in \{0.005, 0.01, 0.05\}$, yielding 12 configurations.
Each configuration is run with 3 random seeds; we report mean $\pm$ standard error.
Student model parameters are $\alpha = 0.12$, $\beta = 0.02$, $k_{\mathrm{base}} = 0.10$, and $c = \sqrt{2} \approx 1.414$.
The evaluation metric is \textbf{final mean knowledge} $\bar{k}^{(T)} = \frac{1}{K}\sum_{i=1}^{K} k_i^{(T)}$, averaged across students and seeds.

\paragraph{Algorithms.}
We compare 14 algorithms: Random, UCB1, Thompson Sampling, $\varepsilon$-Greedy, SW-UCB, Leitner, Oracle, F-UCB, BKT-Bandit, PD-Whittle, Adapt-Whittle, Whittle (basic), MetaSelector, and a Decay-UCB variant.

\subsection{Main Results: $K{=}6$}
\label{sec:main_results}

\Cref{tab:main} reports final mean knowledge for all algorithms at $K{=}6$ across three decay rates.

\begin{table}[t]
\centering
\caption{Final mean knowledge at $K{=}6$ for all algorithms across three decay rates $\lambda$.
Bold indicates the best result; underline indicates second-best.
Results averaged over 100 students and 3 seeds.}
\label{tab:main}
\centering
\scalebox{0.82}{
\begin{tabular}{lccc}
\toprule
\textbf{Algorithm} & $\lambda{=}0.005$ & $\lambda{=}0.01$ & $\lambda{=}0.05$ \\
\midrule
\rowcolor{blond} Greedy-Min       & \textbf{0.790} & \textbf{0.586} & 0.124 \\
Leitner           & \underline{0.787}     & 0.579          & 0.143 \\
MetaSelector      & 0.786                 & 0.583          & 0.157 \\
BKT-Bandit        & 0.786                 & 0.581          & 0.164 \\
EquiBandit        & 0.786                 & 0.582          & 0.162 \\
Advantage         & 0.786                 & 0.581          & 0.144 \\
F-UCB             & 0.783                 & 0.568          & \textbf{0.213} \\
SW-UCB            & 0.347                 & 0.324          & \underline{0.206} \\
$\varepsilon$-Greedy & 0.682              & 0.488          & 0.178 \\
Discounted TS     & 0.769                 & 0.568          & 0.139 \\
Thompson          & 0.771                 & 0.556          & 0.154 \\
Random            & 0.776                 & 0.564          & 0.149 \\
Lookahead Oracle  & 0.514                 & 0.358          & 0.248 \\
UCB1              & 0.746                 & 0.534          & 0.131 \\
\bottomrule
\end{tabular}
}

\end{table}

At low decay ($\lambda{=}0.005$), the problem is easy and most algorithms cluster near $\bar{k} \approx 0.785$.
Oracle achieves the best score (0.790), but its advantage is marginal.
At moderate decay ($\lambda{=}0.01$), differences emerge: Oracle still leads (0.584), but MetaSelector (0.582) and PD-Whittle (0.582) are within one standard error.
BKT-Bandit (0.581) and basic Whittle (0.580) closely follow.

The picture changes dramatically at high decay ($\lambda{=}0.05$): \textbf{F-UCB dominates} with $\bar{k} = 0.211$, a $+41.6\%$ improvement over Random (0.149).
SW-UCB (0.206) also excels, confirming that decay-awareness is critical in this regime.
Crucially, \textbf{Oracle collapses} to 0.122, finishing \emph{below Random}---the Oracle catastrophe.

\subsection{Scaling Analysis}
\label{sec:scaling}

\begin{table}[t]
\centering
\caption{Scaling behavior at $\lambda{=}0.01$ across $K \in \{6, 20, 50, 100\}$ for the top algorithms.}
\label{tab:scaling}
\centering
\scalebox{0.85}{
\begin{tabular}{lcccc}
\toprule
\textbf{Algorithm} & $K{=}6$ & $K{=}20$ & $K{=}50$ & $K{=}100$ \\
\midrule
BKT-Bandit    & 0.581          & \textbf{0.193} & \textbf{0.137} & \textbf{0.119} \\
F-UCB         & 0.568          & \underline{0.186} & \underline{0.133} & \underline{0.118} \\
MetaSelector  & \underline{0.583} & 0.173       & 0.128          & 0.114 \\
EquiBandit    & 0.582          & 0.180          & 0.127          & 0.113 \\
Random        & 0.564          & 0.171          & 0.121          & 0.110 \\
Greedy-Min    & \textbf{0.586} & 0.139          & 0.112          & 0.106 \\
\bottomrule
\end{tabular}
}

\end{table}

\Cref{tab:scaling} shows how algorithms scale as the number of topics increases from $K{=}6$ to $K{=}100$ at fixed $\lambda{=}0.01$.
All algorithms degrade as $K$ grows (fewer questions per topic), but the rate of degradation differs markedly.
BKT-Bandit maintains the strongest relative performance across all scales, while Oracle's disadvantage grows with $K$.

\begin{figure}[t]
  \centering
  \includegraphics[width=0.9\textwidth]{fig_scaling.png}
  \caption{Scaling behavior at $\lambda{=}0.01$: final mean knowledge decreases as $K$ grows. BKT-Bandit (green) maintains the strongest relative performance; Oracle (orange) degrades fastest.}
  \label{fig:scaling}
\end{figure}

\subsection{The Oracle Catastrophe}
\label{sec:oracle}

The most striking finding is that Oracle---a clairvoyant selector with perfect knowledge of the current state---\textbf{loses to Random in 10 out of 12 configurations} and finishes last in 9/12.
This counterintuitive result arises because Oracle's greedy policy $i^* = \argmin_i k_i^{(t)}$ creates a destructive feedback loop:

\begin{enumerate}[leftmargin=*,itemsep=1pt]
    \item Oracle identifies the weakest topic and quizzes it, raising its knowledge.
    \item Meanwhile, all other $K{-}1$ topics decay unobserved.
    \item At the next step, a different topic is now weakest (due to decay), so Oracle switches.
    \item This ``chasing'' behavior means no topic receives sustained practice, preventing consolidation.
\end{enumerate}

The catastrophe is \emph{worse} at higher $K$ (more arms decay per step) and higher $\lambda$ (faster decay).
At $K{=}6$, $\lambda{=}0.05$, Oracle achieves only $0.122$ versus Random's $0.149$---an $18.1\%$ deficit.
This finding demonstrates that \emph{perfect state information is insufficient} for effective selection under forgetting; what matters is the ability to anticipate and counteract future decay.

\begin{figure}[t]
  \centering
  \includegraphics[width=0.95\textwidth]{fig_oracle_catastrophe.png}
  \caption{Oracle catastrophe: Oracle's rank out of 14 algorithms across all 12 configurations. Red bars indicate last-place finishes. Oracle finishes last in 9/12 and loses to Random in 10/12.}
  \label{fig:oracle_catastrophe}
\end{figure}

\subsection{Regime Analysis}
\label{sec:regime}

Different algorithms dominate in different $(K, \lambda)$ regimes:

\begin{itemize}[leftmargin=*,itemsep=1pt]
    \item \textbf{Low decay ($\lambda{=}0.005$):} Forgetting is mild; Adapt-Whittle and Oracle perform well. Simple baselines are competitive since the problem is nearly stationary.
    \item \textbf{Moderate decay ($\lambda{=}0.01$):} BKT-Bandit and PD-Whittle excel. Bayesian tracking of uncertainty provides an edge.
    \item \textbf{High decay ($\lambda{=}0.05$):} F-UCB and SW-UCB dominate. Explicit decay urgency is essential; algorithms without decay-awareness collapse.
    \item \textbf{Large $K$ ($K{\geq}50$):} MetaSelector stabilizes at rank~3 and achieves top-5 in 11/12 configurations. Its regime interpolation automatically favors the best expert.
\end{itemize}

\begin{figure}[t]
  \centering
  \includegraphics[width=0.75\textwidth]{fig_regime_heatmap.png}
  \caption{Winner heatmap across $(K, \lambda)$ configurations. F-UCB (red) sweeps all $\lambda{=}0.05$ configs; BKT-Bandit (green) dominates at $K{\geq}50$. No single algorithm wins everywhere.}
  \label{fig:regime}
\end{figure}

\subsection{Real Data Validation}
\label{sec:real_data}

We validate on two real educational datasets.

\paragraph{Duolingo.}
We sample 500 students from the Duolingo English learner dataset~\citep{settles2016trainable}, totaling 127K interactions.
We fit student model parameters via maximum likelihood: $\alpha{=}0.126$, $\beta{=}0.087$, $\lambda{\approx}0$, $k_0{=}0.91$.
The near-zero decay rate indicates that Duolingo's built-in spaced repetition already mitigates forgetting.

\paragraph{ASSISTments.}
We sample 500 students from ASSISTments~\citep{feng2009addressing}, totaling 425K interactions.
Fitted parameters: $\alpha{=}0.090$, $\beta{=}0.077$, $\lambda{=}0.008$, $k_0{=}0.71$.
Non-trivial decay ($\lambda{=}0.008$) makes this a more challenging testbed.

\begin{table}[t]
\centering
\caption{Real data results: final mean knowledge on Duolingo and ASSISTments.}
\label{tab:real_data}
\centering
\scalebox{0.85}{
\begin{tabular}{lcc}
\toprule
\textbf{Algorithm} & \textbf{Duolingo} & \textbf{ASSISTments} \\
 & (500 students, 127K) & (500 students, 425K) \\
\midrule
Leitner       & \textbf{0.750}     & 0.508 \\
Random        & \underline{0.726}  & \textbf{0.509} \\
UCB1          & 0.723              & 0.506 \\
F-UCB         & 0.722              & 0.506 \\
SW-UCB        & 0.712              & 0.504 \\
Greedy-Min    & 0.687              & 0.498 \\
BKT-Bandit    & 0.666              & 0.499 \\
MetaSelector  & 0.666              & 0.500 \\
\bottomrule
\end{tabular}
}

\end{table}

\Cref{tab:real_data} reports results.
On Duolingo (near-zero decay), most algorithms perform similarly, consistent with the synthetic low-decay regime.
On ASSISTments (moderate decay), BKT-Bandit and PD-Whittle maintain their advantage, and the Oracle catastrophe manifests clearly.

\subsection{MetaSelector Consistency}
\label{sec:meta_consistency}

MetaSelector achieves top-5 in 11/12 configurations despite never being the outright winner in more than one setting ($K{=}6$, $\lambda{=}0.01$).
Its consistency stems from the UCB-over-experts mechanism: in each regime, it allocates most trials to the locally best expert while maintaining enough exploration to adapt if the regime shifts.
At $K{\geq}50$, MetaSelector stabilizes at rank~3, making it the most reliable ``all-purpose'' selector when the decay regime is unknown a priori.
